\chapter{交割、资金划付与投后治理}

\section{交割条件（CP）}

常见 CP 包括：陈述与保证真实、无法律禁令、关键第三方同意、政府审批（若需）、员工留任协议签署、无重大不利影响（MAC）未发生。GP 应维护\textbf{CP 检查表}，由律师出具法律意见、财务确认资金路径。\textbf{划款指令}须双人复核，防范欺诈性邮件篡改账户。

\section{文件签署顺序}

通常先满足 CP，再签署最终股份认购/增资文件，最后释放托管或直接划款。跨境交易可能涉及\textbf{外汇登记}与税务代扣；时间表应预留监管机构处理窗口。签字页与正本份数在 closing memo 中列明。

\section{投后信息权}

股东协议与 LPA 往往规定：月度/季度报表、年度预算、重大事件即时通报、审计配合权。GP 应建立\textbf{被投企业报告模板}与内部 CRM 提醒，避免投后「失联」引发 LP 质疑。

\section{董事会节奏与赋能}

董事会频率与议程应匹配公司阶段：早期可更密，成熟期可季度。GP 董事应在\textbf{战略、资本结构、关键招聘}上发力，避免微观管理日常运营。赋能形式包括介绍客户、补充高管、协助下一轮融资——均须注意利益冲突披露。

\section{重大事项同意权}

基金层面可能对超过单笔金额或特定类型的投资要求 LPAC 或顾问委员会知情；被投企业层面的否决权行使应记录在案，以支持未来退出时的治理叙事。
